
\section{Symbolizing English}

% \subsection{First steps to symbolization} \cite[][Chapter 4, pp. 27-31]{Mag25}
% \subsection{Connectives} \cite[][Chapter 5, pp. 32-48]{Mag25}

Last week we showed how to use capital letters as stand-ins for sentences
in logical arguments and described this is as fundamental to formal logics.
This week we will discuss this process--often called \textit{symbolization}
or \textit{translation}--in more detail. 
\\ \\ 
Recall that we can describe a argument using sentence letters to reveal its
deeper formal structure. For example, we can stipulate that each of the
following arguments is \textit{valid in virtue of its form}.

\begin{argument}
\item[] A
\item[] If A, then C
\item[\therefore] C
\end{argument}

\begin{argument}
\item[] A or B
\item[] not-A
\item[\therefore] B
\end{argument}

\begin{argument}
\item[] not-(A and B)
\item[] A
\item[\therefore] not-B
\end{argument}

The logic of sentences that we define is called Truth-Functional Logic (TFL)
or Propositional Logic (PL), depending on the author. We will use TFL but take
the names to be synomymous references to the same logic.
\\ \\ 
Note that by limiting ourselves to sentence letters, we are marking off an
explicit boundary in our logic between those arguments that depend in some
way on the 'content' of their sentences and those that depend on form alone.
\\ \\
When we look at Predicate Logic (Quantificational Logic) later this semester,
we will see how to incorporate the contributions of names and properties so
that we can encode \textbf{some} of the contributions that (sub-sentential)
content makes to the rules of inference. But, for now, we will focus on
sentences and the form of the arguments that they appear in.

\begin{quote}
  At best, TFL will help us to understand arguments that are valid due to
  their form. \cite[][p30]{Mag25}
\end{quote}

In the examples above, we have some sentences that contain \textit{sentence
letters} themselves. For example: \textit{If A, then B}. When we use capital
letters to stand for sentences (or subsentences) in this way, we are assuming
that the sentences represented are somewhat simple and can be assessed as true
or false. We will call sentences of this sort \textit{atomic sentences}. They
cannot be broken down into parts and we cannot 'peek' inside them to support
our inferences.
\\ \\
When we use letters to represent sentences in this way, we will often provide a
\textit{symbolization key} to help decode the arguments.

\begin{argument}
\item[A:] It is a perfect day for fetch.  
\item[C:] Bandit is happy.
\end{argument}

Of course, these devices are \textit{ad hoc} since TFL is a formal logic and
the actual contents of its sentences is immaterial to the analysis.
\\ \\ 
We can build complex sentences from simpler ones (including atomic sentences)
with the help of a few \textit{connectives}. As the name suggests, the
connectives of TFL are \textit{truth-functional}. In other words, we will
interpret the truth of a complex sentence as a function of the truth of its
component subsentences based on the meanings of its connectives and the order
that they are introduced.
\\ \\
Consider the simple one-place connective $\lnot$. We can understand this
connective as a function from a truth-value to its opposite. So, for example,
when $A$ is true, $\lnot A$ is false. We can summarize the basic TFL
connectives as in the following table.

\begin{table}[h]
\centering
\begin{tabular}{|l|l|l|}
\hline
\textbf{Function} & \textbf{Symbol} & \textbf{Paraphrase} \\
\hline
Negation & $\lnot$ & 'It is not the case that \dots' \\
Conjunction & $\land$ & 'Both \dots and \dots' \\
Disjunction & $\lor$ & 'Either \dots or \dots' \\
Conditional & $\rightarrow$ & 'If \dots then \dots' \\
Biconditional & $\leftrightarrow$ & '\dots if and only if \dots' \\
\hline
\end{tabular}
\caption{Basic TFL Connectives}
\label{tab:tflconnectives}
\end{table}

The symbolization of sentences is subtle and not always a simple preservation
of the surface forms found in the English version. Consider each of the
connectives in turn.

\begin{argument}
\item[R:] A recession is inevitable.
\item[P:] A recession can be prevented.
\end{argument}

If $P$ is part of our symbolization key, we might want to symbolize $R$ as
$\lnot P$ to capture the relationship between the sentences, even though
\textit{not} and \textit{it is not the case that} do not occur anywhere in
the sentence. Here is where the concept of \textit{paraphrase} comes into
play. We will want to lean pretty heavily on the flexibility of paraphrase,
especially when dealing with words that have antonyms or other 'lexical'
negations with \textit{un-}, \textit{non-}, \textit{de-}, and so on. 
\\ \\
We see a similar separation of English forms and TFL forms in our
treatment of conjunctions. If we want to symbolize \textit{Barbara is a medical
  doctor and a philosopher}, we need to depart from the English form and treat
it as a conjunction of sentences ($M \land P$) instead of as one sentence
(\textit{Barbara M and P}, perhaps). That is, we need to mention Barbara in
each of the sentences, not just the once. Our symbolization key would be:

\begin{argument}
\item[M:] Barbara is a medical doctor.
\item[P:] Barbara is a philosopher.  
\end{argument}

Note again that the semantics of TFL is bare bones. We will bundle under
conjunction several English connectives than have slightly different
implications in their full meaning: \textit{and}, \textit{but}, and
\textit{although} in the defintion, but many more in practice. In addition to
collapsing these English connectives to one in TFL, we will also ignore the
implications of sentence order. 

\begin{argument}
\item[A:] John left.
\item[B:] Mary got upset.
\end{argument}

Although \textit{A and B} and \textit{B and A} have different implications
about the order of events--and perhaps even about cause--in TFL $A \land B$
says the exact same thing as $B \land A$. But the interaction of the operators
is important. The following have a different meaning in TFL, even though each
could be the correct symbolization of the English sentence schema \textit{It
  is not the case that A and B}.
\\
\\
\twoitems{$\lnot(A \land B)$}{$(\lnot A \land B)$}
\\
\\
We need to take care in the order we apply the connectives.
\\
\\
With these preliminaries out of the way, we can now provide informal
definitions for each of the basic TFL connectives by leaning on English
as our \textit{metalanguage} and expressing the definitions by way of
paraphrase. We will use lower-case Greek letters $\phi$, $\psi$, $\chi$,
\dots in our metalanguage to represent arbirtary TFL sentences.

\begin{definition}[Negation]
  A sentence can be symbolized as $\lnot \phi$ if it can be paraphrased as
  'It is not the case that $\phi$.'
\end{definition}

\begin{definition}[Conjunction]
  A sentence can be symbolized as $\phi\land\psi$ if it can be paraphrased as
  'Both $\phi$ and $\psi$', as '$\phi$, but $\psi$' or as 'Although $\phi$,
  $\psi$'.
\end{definition}

\begin{definition}[Disjunction]
  A sentence can be symbolized as $\phi\lor\psi$ if it can be paraphrased as
  'Either $\phi$ or $\psi$'.
\end{definition}

\begin{definition}[Conditional]
  A sentence can be symbolized as $\phi\rightarrow\psi$ if it can be
  paraphrased as 'If $\phi$, then $\psi$' or as '$\psi$ only if $\phi$'.
\end{definition}

\begin{definition}[Biconditional]
  A sentence can be symbolized as $\phi\leftrightarrow\psi$ if it can be
  paraphrased as '$\phi$ if and only if $\psi$'.
\end{definition}
